\PilotPage{21}

Consider a vertical spring--mass vibration system consisting of a particle of mass $m$ hanging on an elastic spring of stiffness $k$ as shown in Figure~2.

{\normalsize\bfseries\color{Graphite} 1.1.8.1\quad Static analysis\par}

Initially, the mass is at rest and in a state of equilibrium with the static tension in the spring (Figure~2a). In this equilibrium state, the weight $W=mg$ is balanced by the static tension in the spring $k\delta_{\mathrm{st}}$, where $\delta_{\mathrm{st}}$ is the static displacement of the spring. The free body diagram of the mass in equilibrium condition is shown in Figure~2b. Force-balance analysis in the vertical direction gives
\begin{equation}
    \sum F_y=mg-k\delta_{\mathrm{st}}=0 \qquad\text{(static equilibrium)} \tag{26}
\end{equation}
Solution of Eq.~(26) yields the static displacement $\delta_{\mathrm{st}}$ as
\begin{equation}
    \delta_{\mathrm{st}}=\frac{mg}{k} \qquad\text{(static displacement)} \tag{27}
\end{equation}

{\normalsize\bfseries\color{Graphite} 1.1.8.2\quad Vibration analysis\par}

If the mass is displaced from this equilibrium state and then let free, it will oscillate up and down about the equilibrium position, i.e., it will experience a state of vibration with time-dependent displacement $u(t)$. The free body diagram of the mass in the displaced position gives the sum of forces in the vertical direction as
\begin{equation}
    \sum F_y=mg-k\left[\delta_{\mathrm{st}}+u(t)\right] \qquad\text{(vibration motion)} \tag{28}
\end{equation}
Expansion of Eq.~(28) gives
\begin{equation}
    % SOURCE CORRECTION: The expansion of Eq.~(28) requires a negative $ku(t)$ term.
    \sum F_y=(mg-k\delta_{\mathrm{st}})-ku(t) \qquad\text{(vibration motion)} \tag{29}
\end{equation}
However, $(mg-k\delta_{\mathrm{st}})=0$ as indicated by the equilibrium analysis in Eq.~(26). Hence, the sum of forces acting on the vibrating mass is simply
\begin{equation}
    % SOURCE CORRECTION: Eq.~(30) must retain the negative spring-restoring force.
    \sum F_y=-ku(t) \qquad\text{(sum of forces for vibration analysis)} \tag{30}
\end{equation}
Newton's law of motion (NLM) for the vertical direction is stated as
\begin{equation}
    \sum F_y=m\ddot{u} \qquad\text{(NLM)} \tag{31}
\end{equation}
Substitution of Eq.~(30) into Eq.~(31) gives, upon rearrangement,
\begin{equation}
    m\ddot{u}(t)+ku(t)=0. \tag{32}
\end{equation}
Eq.~(32) is the same as Eq.~(4) developed previously for the horizontal spring--mass vibrating system. Its general solution is that given by Eq.~(21), i.e.,
\begin{equation}
    u(t)=C\sin(\omega_n t+\varphi). \tag{21}
\end{equation}
